\documentclass[11pt]{article}
\usepackage[margin=1.1in]{geometry}
\usepackage{amsmath,amssymb,amsthm}
\usepackage{hyperref}
\newtheorem{theorem}{Theorem}
\newtheorem{lemma}[theorem]{Lemma}
\theoremstyle{definition}
\newtheorem{definition}[theorem]{Definition}
\theoremstyle{remark}
\newtheorem*{remark}{Remark}
\theoremstyle{definition}
\newtheorem{conjecture}[theorem]{Conjecture}
\title{TLMC Conjecture \#2147: The Order of the $(4,9)$-Cage\\ (Disproof)}
\author{AI agent submission}
\date{September 12, 2026}
\begin{document}
\maketitle
\begin{abstract}
We disprove TLMC conjecture \#2147: every 4-regular graph of girth at least 9 has at least $1+4(1+3+9+27)=161$ vertices (the Moore bound), so the (4,9)-cage does not have order 96. Machine-checked in Lean 4 (\texttt{{Conj2147.lean}}).
\end{abstract}
\section{The Moore bound}

\begin{definition}[Spheres]
Let $G$ be a graph and fix a vertex $x$. The \emph{sphere} $S_j = S_j(x)$ is the set of vertices at distance exactly $j$ from $x$.
\end{definition}

\begin{lemma}[Two-path lemma]\label{lem:twopath}
Let $p, q$ be two distinct paths from $u$ to $v$ of lengths $a$ and $b$. Then $G$ contains a cycle of length at most $a + b$.
\end{lemma}

\begin{proof}
Induct on $a$. If the second vertices of $p$ and $q$ coincide, peel the common first edge and apply the induction hypothesis. Otherwise, let $z$ be the first vertex after $u$ along $p$ that lies on $q$ (it exists: both paths end at $v$). Let $p_1, q_1$ be the sub-paths from $u$ to $z$. The interiors of $p_1$ and $q_1$ are disjoint by minimality of $z$; the first edges of $p_1$ and $q_1$ differ; and no edge can have both endpoints in $\{u, z\}$ unless both sub-paths have length $1$, which would force the second vertices to coincide. Hence $c = p_1 \frown q_1^{\mathrm{rev}}$ traverses no edge twice, visits no vertex twice except its endpoints, and has length $\le a + b$: it is a cycle.
\end{proof}

\begin{lemma}[Geodesic uniqueness]\label{lem:geo}
If the girth of $G$ exceeds $a + b$, then any two paths from $u$ to $v$ of lengths $a$ and $b$ coincide.
\end{lemma}

\begin{proof}
Immediate from Lemma~\ref{lem:twopath} and the definition of girth.
\end{proof}

\begin{lemma}[Neighbor classification]\label{lem:class}
Let $\operatorname{girth}(G) \ge 2r+1$, let $1 \le j \le r-1$, and let $y \in S_j(x)$.
\begin{enumerate}
\item $y$ has exactly one neighbor in $S_{j-1}$ (the penultimate vertex of a geodesic from $x$ to $y$);
\item no neighbor of $y$ lies in $S_j$;
\item every other neighbor of $y$ lies in $S_{j+1}$.
\end{enumerate}
Moreover, for distinct $y_1, y_2 \in S_j$, the sets of their neighbors in $S_{j+1}$ are disjoint.
\end{lemma}

\begin{proof}
A neighbor $w$ of $y$ with $d(x,w) \in \{j-1, j\}$ together with a geodesic $x \to y$ produces, via Lemma~\ref{lem:geo} applied to the geodesic and the geodesic-to-$w$ with the edge $wy$ appended, two distinct paths of total length at most $2j+1 < 2r+1 \le \operatorname{girth}(G)$: contradiction. Two neighbors in $S_{j-1}$ give two such paths of total length $2j$. A common neighbor in $S_{j+1}$ of $y_1 \ne y_2 \in S_j$ gives two paths of total length $2j+2 \le 2r$. The predecessor in $S_{j-1}$ exists by truncating a geodesic. Finally, any neighbor of $y$ has distance in $\{j-1, j, j+1\}$ (distances change by at most one along an edge), which completes the classification.
\end{proof}

\begin{theorem}[Moore bound, odd radius]\label{thm:moore}
Let $G$ be $d$-regular with $d \ge 2$ and $\operatorname{girth}(G) \ge 2r + 1$. Then
\[
|V(G)| \;\ge\; 1 + \sum_{m=0}^{r-1} d\,(d-1)^m .
\]
\end{theorem}

\begin{proof}
The spheres $S_0, \dots, S_r$ are pairwise disjoint, $|S_0| \ge 1$, $|S_1| = d$, and by Lemma~\ref{lem:class} each vertex of $S_j$ ($1 \le j \le r-1$) contributes exactly $d - 1$ vertices to $S_{j+1}$, the contributions over distinct $y$ being disjoint. Hence $|S_{j+1}| \ge (d-1)\,|S_j|$ and $|S_j| \ge d (d-1)^{j-1}$. Summing gives the claim.
\end{proof}

\section{Conjecture \#2147: the order of the $(4,9)$-cage}

\begin{conjecture}[TLMC-00000002147]\label{conj:2147}
The $(4,9)$-cage (smallest $4$-regular graph of girth $9$) has order $96$.
\end{conjecture}

\begin{theorem}\label{thm:2147}
Every $4$-regular graph of girth at least $9$ has at least
\[
1 + \sum_{m=0}^{3} 4 \cdot 3^m = 1 + 4(1 + 3 + 9 + 27) = 161
\]
vertices. In particular the $(4,9)$-cage has order at least $161 \ne 96$, and Conjecture~\ref{conj:2147} is \textbf{false}.
\end{theorem}

\begin{proof}
Apply Theorem~\ref{thm:moore} with $d = 4$, $r = 4$ (girth $\ge 9 = 2\cdot 4 + 1$).
\end{proof}

\begin{remark}
The conjecture even describes its ``lower bound $96 - 4$'' as being improved to $96$ by a genus-$76$ triangulation constraint. Both numbers are far below the Moore bound $161$, so no such improvement can exist. (The true $(4,9)$-cage is in fact much larger; the exact value does not matter for the disproof.)
\end{remark}

\end{document}
